% !TeX root = ../../Book1.tex
% 纲要标题：### 21.1 分圆多项式
% 纲要位置：计划纲要.md，第 705 行。

\section{分圆多项式}
\label{b1:sec:2101}

上一章在“基域已经含有所需单位根”的条件下讨论开方。现在把基域取为 \(\mathbb Q\)，研究加入单位根本身造成的扩张。需要区分两件事：\(X^n-1\) 收集全部 \(n\) 次单位根，而分圆多项式只收集乘法阶恰为 \(n\) 的那些根。

% 纲要内容（仅作写作计划，尚非教材正文）：
%
% * 单位根与分圆多项式
% * 分圆多项式的不可约性
% * 分圆域及其 Galois 群
%

\subsection{单位根与分圆多项式}
\label{b1:subsec:2101:01}

取 \(\zeta_n=e^{2\pi i/n}\)。它是一个本原 \(n\) 次单位根；所有本原 \(n\) 次单位根为 \(\zeta_n^a\)，其中 \(1\leq a\leq n\)、\(\gcd(a,n)=1\)。

\begin{definition}\label{b1:def:cyclotomic-polynomial}
设 \(n\geq1\) 为整数，并在 \(\mathbb C\) 中选取一个本原 \(n\) 次单位根 \(\zeta_n\)。多项式
\[
\Phi_n(X)\coloneqq\prod_{\substack{1\leq a\leq n\\\gcd(a,n)=1}}(X-\zeta_n^a)
\]
不依赖本原单位根 \(\zeta_n\) 的选择，称为第 \(n\) 个\term[fenyuan-duoxiangshi]{分圆多项式}[cyclotomic polynomial]。其次数为 Euler 函数 \(\varphi(n)\)，即模 \(n\) 可逆剩余类的个数。
\end{definition}
\symidx{\(\Phi_n(X)\)：第 \(n\) 个分圆多项式}

每个 \(n\) 次单位根的阶是 \(n\) 的唯一一个正因子，因此
\begin{equation}\label{b1:eq:cyclotomic-factorization}
X^n-1=\prod_{d\mid n}\Phi_d(X).
\end{equation}
这还给出 \(\Phi_n\in\mathbb Z[X]\)。对 \(n\) 归纳，\(\Phi_1=X-1\)；若真因子对应的多项式均已知为首一整系数多项式，则用它们的乘积去除 \(X^n-1\)。首一除法的商、余式仍在 \(\mathbb Z[X]\) 中，而\eqref{b1:eq:cyclotomic-factorization}说明在 \(\mathbb C[X]\) 中余式为零、商为 \(\Phi_n\)，故结论成立。

\ProseParagraphBreak
例如 \(\Phi_2=X+1\)、\(\Phi_3=X^2+X+1\)、\(\Phi_4=X^2+1\)、\(\Phi_6=X^2-X+1\)。对素数 \(p\)，有 \(\Phi_p=1+X+\cdots+X^{p-1}\)。不过，整系数性本身不能说明不可约；下一小节需要另外论证所有本原根互为有理数域上的共轭。

\subsection{分圆多项式的不可约性}
\label{b1:subsec:2101:02}

\begin{theorem}\label{b1:thm:cyclotomic-irreducible}
设 \(n\geq1\) 为整数，\(\zeta_n\in\mathbb C\) 为本原 \(n\) 次单位根。则第 \(n\) 个分圆多项式 \(\Phi_n(X)\) 在 \(\mathbb Q[X]\) 中不可约。因此 \(\Bigl[\mathbb Q(\zeta_n):\mathbb Q\Bigr]=\varphi(n)\)。
\end{theorem}
\begin{proof}
设 \(f\) 为 \(\zeta_n\) 在 \(\mathbb Q\) 上的首一最小多项式。它是首一整系数多项式 \(\Phi_n\) 的因子，由\cref{b1:thm:gauss-descent}，可写 \(\Phi_n=fg\)，其中 \(f,g\in\mathbb Z[X]\) 都首一。

先证明：若 \(\alpha\) 是 \(f\) 的根，\(q\) 为不整除 \(n\) 的素数，则 \(\alpha^q\) 仍为 \(f\) 的根。否则 \(\alpha^q\) 是 \(\Phi_n\) 的根而不是 \(f\) 的根，所以是 \(g\) 的根；从 \(g(\alpha^q)=0\) 得 \(f(X)\mid g(X^q)\) 于 \(\mathbb Q[X]\)。首一除法又说明商在 \(\mathbb Z[X]\) 中。

把这个整系数整除关系模 \(q\) 约化，并以横线表示所得多项式。特征 \(q\) 的二项式公式给出
\[
\bar f(X)\mid\bar g(X^q)=\bar g(X)^q.
\]
\(\bar f\) 首一且次数为正，故有一个不可约因子 \(h\)。整除关系迫使 \(h\mid\bar g\)，因而 \(h^2\mid\bar f\bar g=\overline{\Phi_n}\)，进而 \(h^2\mid X^n-1\) 于 \(\mathbb F_q[X]\)。但是 \(q\nmid n\)，所以 \(X^n-1\) 与导数 \(nX^{n-1}\) 互素，不可能有重不可约因子，矛盾。

任意与 \(n\) 互素的正整数 \(a\) 都是若干不整除 \(n\) 的素数的乘积。反复应用刚证的断言，得到 \(\zeta_n^a\) 是 \(f\) 的根。因此 \(f\) 含有全部 \(\varphi(n)\) 个本原根；而 \(f\mid\Phi_n\)，只能有 \(f=\Phi_n\)。次数结论来自\cref{b1:prop:minimal-polynomial-degree}。
\end{proof}

\subsection{分圆域及其 Galois 群}
\label{b1:subsec:2101:03}

\begin{theorem}\label{b1:thm:cyclotomic-galois-group}
设 \(n\geq1\) 为整数，\(\zeta_n\in\mathbb C\) 为本原 \(n\) 次单位根。则\term[fenyuan-yu]{分圆域}[cyclotomic field] \(\mathbb Q(\zeta_n)\) 是 \(X^n-1\) 在 \(\mathbb Q\) 上的分裂域，且
\[
\operatorname{Gal}\Bigl(\mathbb Q(\zeta_n)/\mathbb Q\Bigr)
\cong(\mathbb Z/n\mathbb Z)^\times,
\qquad \sigma_a(\zeta_n)=\zeta_n^a.
\]
\end{theorem}
\begin{proof}
全部 \(n\) 次单位根都是 \(\zeta_n\) 的幂，故域正规；特征为零保证可分。自同构保持元素乘法阶，因此必须把 \(\zeta_n\) 送到某个本原根 \(\zeta_n^a\)，这给出到单位群的单射。由\cref{b1:thm:cyclotomic-irreducible}，两端阶均为 \(\varphi(n)\)，故满射。复合在生成元上的作用是 \(\zeta_n\mapsto\zeta_n^{ab}\)，所以这是群同构。
\end{proof}

例如 \(\mathbb Q(\zeta_5)/\mathbb Q\) 的群为 \(C_4\)，而 \(\mathbb Q(\zeta_8)/\mathbb Q\) 的群为 \(C_2\times C_2\)，不能把所有分圆扩张都称为循环扩张。复共轭总对应于剩余类 \(-1\)；当 \(n\geq3\) 时，其固定域为 \(\mathbb Q(\zeta_n+\zeta_n^{-1})\)。确实，后一个域由实数生成，且 \(\zeta_n\) 满足
\[
X^2-(\zeta_n+\zeta_n^{-1})X+1=0;
\]
\(\zeta_n\) 本身非实，因此这一步的次数恰为 \(2\)，与共轭子群的固定域次数相同。
